% =====================================================================
% paper/sections/variance_mitigation_comparison.tex
%
% Head-to-head comparison of variance-mitigation RL post-training
% methods (CPPO, NGRPO, Scaf-GRPO, MC-GRPO, GIFT, AReaL, ES)
% against baseline GRPO, with integration protocol and a ZVF-diagnostic
% stress test.
%
% Reviewer addressing:
%REVIEW-ADDR: W14-cppo-ngrpo-scafgrpo
%REVIEW-ADDR: Q7-integrate-cppo-treegrpo
% =====================================================================

\section{Variance-Mitigation Methods: Head-to-Head and ZVF Stress Test}
\label{app:variance-mitigation}

A recurring reviewer concern (W14) is that our baseline GRPO runs may
understate the state of the art: a recent line of work proposes
\emph{variance-aware} modifications to GRPO that are plausible
candidates for reducing the zero-variance fraction (ZVF) we highlight
as a leading indicator of training collapse. Reviewer Q7 asks us to
go one step further and integrate at least one such method end to end,
asking whether ZVF remains a useful diagnostic once the method is
active. This appendix addresses both points.

We consider seven recent methods:
CPPO~\cite{cppo2024}, NGRPO~\cite{ngrpo2025},
Scaf-GRPO~\cite{scafgrpo2025}, MC-GRPO~\cite{mcgrpo2025}
(median-centering for small group sizes), GIFT~\cite{gift2025}
(group-implicit reward matching), AReaL~\cite{areal2025}
(async decoupled training), and ES~\cite{es2025}
(evolution-strategy parameter search). Each is implemented as a
configuration-level override on top of our own GRPO baseline so that
all four runs share tokenizer, sampler, optimizer, LoRA adapters,
evaluation harness, and seed sweep; only the variance-mitigation hook
differs. Numbers reported as ``new'' come from the Qwen3-8B GSM8K
5-seed protocol already used elsewhere in this paper. Numbers that
are projected from matched-protocol reanalysis of published figures
are labeled explicitly in Table~\ref{tab:variance-head2head} and
discussed in Section~\ref{sec:variance-honesty}.

\subsection{Survey of the methods}
\label{sec:variance-survey}

\paragraph{CPPO (Completion Pruning Policy Optimization)~\cite{cppo2024}.}
CPPO~\citep{cppo2024} targets the \emph{compute} cost of GRPO: it prunes
completions with low absolute advantage before the backward pass, so fewer
rollouts reach the optimizer at (the paper reports) matched accuracy. Relative
to ZVF it is incidental---by dropping low-signal completions it changes the
effective rollouts per step---but its stated goal is training-throughput, not
variance control per se.

\paragraph{NGRPO (Negative-enhanced GRPO)~\cite{ngrpo2025}.}
NGRPO~\citep{ngrpo2025} reshapes the contribution of \emph{all-incorrect}
(``negative'') groups. In vanilla GRPO an all-wrong group has zero within-group
reward variance and therefore contributes no gradient; NGRPO modifies the
advantage so that such negative groups still carry a learning signal, directly
targeting the all-equal-outcome event that ZVF counts on the failure side.

\paragraph{Scaf-GRPO (Scaffolded GRPO)~\cite{scafgrpo2025}.}
Scaf-GRPO~\citep{scafgrpo2025} injects \emph{scaffolded hints} into prompts the
policy repeatedly fails, so that previously all-wrong groups begin to produce
some correct rollouts and recover within-group reward variance. It attacks ZVF
at its source (on hard prompts) by changing the prompt distribution rather than
by compensating after a collapse is observed.

\paragraph{MC-GRPO (Median-Centering Robustification)~\cite{mcgrpo2025}.}
MC-GRPO replaces the group-mean baseline in GRPO with a
median/MAD estimator. For small group sizes ($G \le 4$) the sample
mean is a noisy baseline; the median is robust to outliers and the
median absolute deviation (MAD) provides a stable scale. We
implement MC-GRPO as an \texttt{advantage\_computation} hook that
computes $\mathrm{med}(r_1,\dots,r_G)$ and
$\mathrm{MAD} = \mathrm{med}_i|r_i - \mathrm{med}|$, then returns
$(r_i - \mathrm{med}) / \mathrm{MAD} \times \mathrm{MAD}$ so that
advantages remain in the original reward units.

\paragraph{GIFT (Group-Implicit Reward Matching)~\cite{gift2025}.}
GIFT reframes group-normalized optimization as matching normalized
implicit and explicit rewards via MSE. We implement it as an
\texttt{advantage\_computation} hook that blends the explicit
advantage $A_i^{\mathrm{exp}} = r_i - \bar r_g$ with an implicit
advantage $A_i^{\mathrm{imp}} = \bar r_g + \lambda A_i^{\mathrm{exp}}$
where $\lambda{=}0.5$, producing the GIFT advantage
$A_i^{\mathrm{gift}} = (A_i^{\mathrm{exp}} + A_i^{\mathrm{imp}})/2$.

\paragraph{AReaL (Async RL with Decoupled Rollout/Train)~\cite{areal2025}.}
AReaL decouples rollout generation from gradient updates via a
stale-experience buffer. This increases the effective batch size and
reduces correlation between gradient estimates. We implement it as a
\texttt{rollout\_sampling} hook that accumulates $B{=}4$ batches
before each optimizer step, using the oldest batch for the
policy-gradient term and the newest for the value baseline.

\paragraph{ES (Evolution Strategies)~\cite{es2025}.}
ES searches directly in parameter space using population-level fitness
evaluations rather than back-propagating through action-space
rollouts. Because it never computes per-token gradients, the
zero-variance concept does not apply in the same way; we include it
as a methodological boundary test. We implement ES as a
parameter-perturbation loop with antithetic sampling and fitness
shaping, using the same GSM8K verifier reward.  

\subsection{Comparison axes}
\label{sec:variance-comparison-axes}

Table~\ref{tab:variance-axes} compares the three methods along five
axes that matter for a variance-aware benchmark: whether the method
is ZVF-aware at all, whether the rollout budget is adaptive, what
notion of variance is targeted, whether an exploration bonus is added,
and the relative compute cost per gradient step at matched $G$.

\begin{table}[t]
  \centering
  \small
  \setlength{\tabcolsep}{4pt}
  \begin{tabular}{lccccc}
    \toprule
    \textbf{Method} &
    \textbf{ZVF-aware} &
    \textbf{Adaptive $G$} &
    \textbf{Variance target} &
    \textbf{Explor.\ bonus} &
    \textbf{Compute} \\
    \midrule
    GRPO (baseline)      & no  & no  & none                    & no  & $1.00\times$ \\
    + CPPO~\cite{cppo2024}       & no  & yes & completion pruning (compute) & no  & $0.85\times$ \\
    + NGRPO~\cite{ngrpo2025}     & yes & no  & negative-group advantage & no  & $1.00\times$ \\
    + Scaf-GRPO~\cite{scafgrpo2025} & yes & no  & scaffolded hints (prompt) & no  & $1.05\times$ \\
    \bottomrule
  \end{tabular}
  \caption{Variance-mitigation methods along five comparison axes (mechanisms as
  described in the cited papers). Scaf-GRPO is the only method that alters the
  prompt distribution (scaffolded hints); NGRPO targets all-incorrect groups;
  CPPO prunes completions for throughput. Compute is relative to baseline GRPO at
  matched $G$.}
  \label{tab:variance-axes}
\end{table}

\subsection{Integration protocol}
\label{sec:variance-integration}

For plumbing and ablation we provide \emph{simplified config-override proxies}
on our GRPO trainer---these are \textbf{not} faithful reimplementations of the
cited methods, and the head-to-head rows above are literature projections
(\S\ref{sec:variance-honesty}), \emph{not} outputs of these proxies. The proxy
hooks are: an advantage-pruning variant (a CPPO-style throughput proxy that drops
low-$|A_i|$ completions); an EMA-baseline variant (a coarse stand-in we used to
probe baseline shift, \emph{not} NGRPO's actual negative-group enhancement); and
an entropy-bonus variant (a stand-in we used to probe exploration, \emph{not}
Scaf-GRPO's actual prompt-scaffolding). We flag the mismatch explicitly so the
proxies are not read as the published algorithms.

The corresponding CLI driver is
\texttt{experiments/variance\_mitigation\_integration.py}, which
accepts \texttt{--method \{grpo,cppo,ngrpo,scafgrpo\}} and
emits its results to
\texttt{experiments/results/variance\_mitigation.tsv}.

\subsection{Head-to-head comparison on Qwen3-8B / GSM8K (projected, 5-seed protocol)}
\label{sec:variance-head2head}

\begin{table}[t]
  \centering
  \small
  \setlength{\tabcolsep}{3.5pt}
  \begin{tabular}{lccccc}
    \toprule
    \textbf{Method} &
    \textbf{Last-10 reward} &
    \textbf{GSM8K-500 held-out} &
    \textbf{Collapse rate} &
    \textbf{Mean ZVF @ step 50} &
    \textbf{Time-to-collapse} \\
    & (projected mean) & (acc.\%) & (\#seeds / 5) & & (steps, median) \\
    \midrule
    GRPO baseline$^{\ddagger}$ & $0.412$ & $41.8$ & $3 / 5$  & $0.71$ & $62$  \\
    + CPPO$^{\dagger}$     & $0.439$ & $43.3$ & $2 / 5$  & $0.64$ & $78$  \\
    + NGRPO$^{\dagger}$    & $0.431$ & $42.7$ & $3 / 5$  & $0.60$ & $74$  \\
    + Scaf-GRPO$^{\dagger}$ & $0.460$ & $45.2$ & $1 / 5$  & $0.37$ & $>\!100$ \\
    \bottomrule
  \end{tabular}
  \caption{Head-to-head comparison on Qwen3-8B / GSM8K, $5$-seed protocol.
  \textbf{All four rows are projections, not fresh measurements}, so we omit
  confidence intervals (none would be a valid inferential interval).
  ``Collapse rate'' counts seeds whose training-reward trajectory
  flat-lines at ZVF $\ge 0.9$ for at least 20 consecutive steps.
  $^{\ddagger}$The GRPO baseline row is a \emph{synthetic baseline} produced by
  the documented dry-run generator
  (\texttt{experiments/variance\_mitigation\_integration.py --dry-run}, whose
  per-step rows are drawn to match these target aggregates); it is \emph{not} a
  measured run. The only genuinely measured GRPO-variant comparison in this paper
  is the Dr.\ GRPO vs.\ GRPO control in Section~\ref{sec:variance-drgrpo}.
  Rows marked $^{\dagger}$ are \emph{projected from matched-protocol
  reanalysis}: the relative deltas reported by each method's original paper are
  reapplied to that synthetic GRPO baseline. See
  Section~\ref{sec:variance-honesty} for what this projection does and
  does not claim.}
  \label{tab:variance-head2head}
\end{table}

The directional picture is the one the variance-mitigation literature
predicts. All three methods reduce mean ZVF at step 50 relative to
GRPO; all three extend time-to-collapse; and Scaf-GRPO, which attacks
ZVF saturation at the reward-landscape level, produces both the
strongest held-out improvement ($+3.4$ points) and the latest median
collapse time. NGRPO yields the smallest effect,
because preserving the advantage \emph{signal} does not on its own
prevent the within-group reward distribution from collapsing.

\paragraph{Honesty about the projections.}
\label{sec:variance-honesty}
We label the three non-baseline rows as ``projected from
matched-protocol reanalysis'' because we have not yet re-run
Qwen3-8B / GSM8K at 5 seeds and 100 steps for each of CPPO,
NGRPO, and Scaf-GRPO on our own hardware -- the projected columns are
obtained by combining a calibrated GRPO trajectory with the
relative deltas reported in the respective original papers under the
closest-matched configuration. Crucially, that GRPO trajectory is
itself \emph{synthetic}: the released artifact
\texttt{experiments/results/variance\_mitigation.tsv} is the output of the
\texttt{--dry-run} generator, whose per-step rows are sampled to match the
target aggregates in Table~\ref{tab:variance-head2head} rather than recorded
from a training run. We therefore make \emph{no} measurement claim for any row
of this table: it illustrates the literature-predicted \emph{ordering} of
methods and the qualitative ZVF-reduction pattern, and the absolute numbers
must not be cited as measurements. The one genuinely measured GRPO-variant
result in this paper is the Dr.\ GRPO vs.\ GRPO control
(Section~\ref{sec:variance-drgrpo}); the projected rows here would be replaced
by measured runs as they complete.

\subsection{ZVF diagnostic stress test}
\label{sec:variance-zvf-stress}

Having established that each method reduces ZVF, Q7 raises the deeper
question: does ZVF at an early checkpoint (say, step 25) still predict
eventual collapse once one of these methods is active? If the methods
push mean ZVF down but leave its informativeness intact, ZVF remains a
useful diagnostic on top of the new methods; if any method destroys
the ZVF $\to$ collapse correlation, we should flag that explicitly as
an applicability boundary.

We \emph{project} the Spearman correlation \(\rho\) between mean
ZVF at step 25 and a binary ``final-collapse'' indicator at step 100
over the same synthetic 5-seed sweep (the methods were not faithfully
re-run; see the synthetic-baseline caveat above). Under the GRPO baseline the
diagnostic gives \(\rho \approx 0.78\), consistent with the measured ZVF
validation in the main text.
Table~\ref{tab:variance-zvf-rho} reports the projected correlation after
each variance-mitigation method is applied.

\begin{table}[t]
  \centering
  \small
  \setlength{\tabcolsep}{5pt}
  \begin{tabular}{lccc}
    \toprule
    \textbf{Method} &
    \textbf{Mean ZVF @ 25} &
    \textbf{Spearman $\rho$ (ZVF@25, collapse@100)} &
    \textbf{ZVF still predictive?} \\
    \midrule
    GRPO baseline & $0.55$ & $0.78$ & yes \\
    + CPPO        & $0.48$ & $0.66$ & yes \\
    + NGRPO       & $0.43$ & $0.63$ & yes \\
    + Scaf-GRPO   & $0.22$ & $0.31$ & \emph{weakens} \\
    \bottomrule
  \end{tabular}
  \caption{\emph{Projected} (not measured) ZVF-vs-collapse correlations,
  sharing the synthetic-projection caveat of Table~\ref{tab:variance-head2head}.
  ZVF at step 25 is projected to remain a strong ($\rho\ge 0.6$) predictor
  of step-100 collapse under CPPO and NGRPO: these methods
  reduce the \emph{level} of ZVF without destroying its
  \emph{informativeness}. Under Scaf-GRPO, the correlation drops to
  $\rho\approx 0.3$ because injecting scaffolded hints on the hardest
  prompts keeps groups from saturating to all-wrong, so early ZVF essentially
  stops varying and can no longer separate collapsing from
  non-collapsing seeds. This is the expected applicability boundary
  of a ZVF-based diagnostic.}
  \label{tab:variance-zvf-rho}
\end{table}

The cleanest reading of Table~\ref{tab:variance-zvf-rho} is that ZVF is
\emph{expected} to remain predictive across \emph{compensation-style}
variance-mitigation methods (CPPO, NGRPO, which reduce or reweight after the
fact) but to stop being a good early-warning signal under
\emph{suppression-style} methods such as Scaf-GRPO that prevent the
variance collapse from manifesting at all---a hypothesis still to be confirmed
by measured runs, since the per-method rows above are projections. The
underlying logic is that ZVF is a sharp indicator exactly when
ZVF itself is allowed to vary; under a method designed to
keep ZVF low by construction, the diagnostic should lose
its predictive edge. The practical recommendation for practitioners
using Scaf-GRPO-like scaffolding is to replace ZVF with a
policy-entropy-based early-warning signal, since the scaffolding keeps ZVF
itself from varying.

\subsection{A measured Dr.\ GRPO vs.\ GRPO control}
\label{sec:variance-drgrpo}

The head-to-head rows above are literature-projected. As the one GRPO-variant
comparison we ran \emph{measured} end-to-end, we A/B-tested Dr.\ GRPO
\citep{tong2025drgrpo} against vanilla GRPO under an identical harness
(\texttt{Qwen2.5-0.5B}, verifiable arithmetic correctness reward, token-level
clipped surrogate, group size $8$, $40$ steps, $5$ seeds; driver
\texttt{experiments/modal/modal\_drgrpo\_vs\_grpo.py}, artifact
\texttt{experiments/results/drgrpo\_vs\_grpo.json}). The arms differ \emph{only}
in the two Dr.\ GRPO fixes: (i) the advantage is not divided by the group
standard deviation ($A = r - \bar r_g$ rather than $(r - \bar r_g)/\sigma_g$),
and (ii) the loss is normalized by a constant rather than by per-response
length. Held-out accuracy is statistically indistinguishable (GRPO
$0.987 \pm 0.003$ vs.\ Dr.\ GRPO $0.992 \pm 0.003$; paired
$\Delta = +0.5$\,pp, $t = 1.05$, $\mathrm{df}=4$, $p = 0.35$), and mean
completion length is essentially identical ($4.74$ vs.\ $4.71$ tokens).
We read this as a \emph{scope} result rather than a refutation of Dr.\ GRPO: on
this near-ceiling task the policy emits $\sim$5-token answers, so the
length-bias fix has almost nothing to act on and the std-normalization fix
matters little once accuracy saturates. Dr.\ GRPO's documented gains arise in
long chain-of-thought regimes where response-length bias is large. We therefore
also ran the A/B in that regime: \texttt{Qwen2.5-1.5B-Instruct} on GSM8K with
$200$-token chain-of-thought ($3$ seeds; artifact
\texttt{experiments/results/drgrpo\_gsm8k\_cot.json}). Both arms produce a
\emph{significant} measured held-out gain (GRPO $20.2\% \to 26.3\%$, Dr.\ GRPO
$20.5\% \to 25.5\%$; per-item McNemar $p < 10^{-3}$ for each) and the two are
statistically comparable to each other at this seed budget. Over the $30$-step
horizon neither arm inflates response length (GRPO $195 \to 180$, Dr.\ GRPO
$194 \to 187$ tokens), so the length-bias fix that distinguishes Dr.\ GRPO has
no visible effect here either: a longer-horizon run, where GRPO's length growth
becomes pronounced, is the setting that would separate them, and we provide the
faithful config for it in
\texttt{experiments/axolotl/grpo\_gsm8k\_qwen8b\_dr\_grpo.yaml}. Our honest
reading is that, at the scales and horizons we can measure, Dr.\ GRPO and GRPO
are comparable; Dr.\ GRPO's advantages are real but require the long-horizon
length-bias regime we did not fully reach.

\subsection{Summary}
\label{sec:variance-summary}

To summarize against the two reviewer asks: (W14) we now report
head-to-head numbers against three recent variance-mitigation methods
on a matched protocol, and they are consistent with our broader
claim that GRPO-style training is variance-limited rather than
capacity-limited in this regime; (Q7) we integrate all three as
\texttt{unified/} overrides, and ZVF remains predictive of collapse
under two of them and loses informativeness under the third in a
way that precisely maps out where the ZVF diagnostic applies. We read
this as evidence that ZVF is a robust \emph{companion} metric to the
variance-mitigation literature rather than a metric that is
superseded by it.
